% discussion.tex -- Enterprise Node Orchestrator Paper (RU-V5)

\section{Discussion}

\subsection{Hypothesis Evaluation}

The hypothesis stated that a Node.js event-driven orchestrator can process
a 64-transaction batch in under 90 seconds with zero data leakage and crash
recovery. Stage~1 results provide partial confirmation:

\begin{itemize}
\item \textbf{Orchestration overhead}: 4,704~ms for batch 64, well below
  the 30-second budget (H0 rejected with 6.4$\times$ margin).

\item \textbf{E2E latency with rapidsnark}: 18.9~s (sequential) or 14.6~s
  (pipelined), meeting the 90~s target with 4.8--6.2$\times$ margin.

\item \textbf{E2E latency with snarkjs}: 156.9~s at batch 64, exceeding
  the target by 1.7$\times$. H0 cannot be rejected for the WASM backend
  at this batch size.

\item \textbf{State machine correctness}: 46/46 tests pass, pipelined
  model validated.

\item \textbf{Crash recovery}: Design supports zero-loss recovery via
  WAL + checkpoint; empirical validation deferred to Stage~3.

\item \textbf{Privacy}: Zero data leakage by construction (node receives
  only hashes); formal audit deferred to Stage~3.
\end{itemize}

\subsection{The Proving Bottleneck}

Across all configurations, ZK proof generation dominates end-to-end latency,
consuming 64\% (rapidsnark b64) to 85\% (snarkjs b8) of the total cycle
time. Orchestration overhead---the sum of batch formation, witness
generation, state management, and scheduling---is a negligible fraction
of the proving budget. This finding is consistent with the push0
architecture~\cite{push0_2025orchestration}, which reports orchestration
overhead of 0.05--0.1\% of proving time.

The practical implication is that node performance optimization should
focus exclusively on prover selection and circuit optimization, not on
orchestration layer engineering. A 2$\times$ improvement in orchestration
overhead (from 593~ms to 296~ms) saves 297~ms, while a 2$\times$
improvement in proving (from 12,757~ms to 6,378~ms) saves 6,379~ms---a
21$\times$ greater impact.

\subsection{Prover Backend Constraint}

The snarkjs WASM backend, while sufficient for batch sizes up to 16,
cannot achieve the 90-second target at batch 64. Three mitigation
strategies are available, ordered by recommendation:

\begin{enumerate}
\item \textbf{Rapidsnark} (recommended for production): C++ native prover
  achieving 4--10$\times$ speedup over snarkjs, reducing batch-64 proving
  from $\sim$150~s to $\sim$12~s. Requires native binary installation.

\item \textbf{Batch size reduction} (safest for MVP): Reducing batch size
  to 16 keeps snarkjs within 31.4~s E2E, at the cost of lower throughput
  (0.53~tx/s vs. 4.37~tx/s).

\item \textbf{GPU prover} (future): ICICLE-Snark~\cite{icicle2025snark}
  achieves 63$\times$ MSM improvement, potentially reducing batch-64
  proving to 1--3~s. Requires GPU infrastructure.
\end{enumerate}

\subsection{Pipeline Speedup Characteristics}

The 1.29$\times$ speedup at batch 64 with rapidsnark is modest but
meaningful: it reduces 10-batch processing from 189~s to 146~s, saving
43~s. The speedup is inherently bounded because the proving phase
dominates---even with perfect overlap, preparation cannot be fully hidden
behind proving when $T_\text{prep} < T_\text{prove+submit}$.

The speedup would increase under two conditions: (1) faster provers that
reduce $T_\text{prove+submit}$ closer to $T_\text{prep}$, or (2) larger
batch sizes that increase $T_\text{prep}$ via witness generation. With
GPU provers (1--3~s proving), preparation time (4.7~s for b64) would
exceed proving time, making the preparation phase the bottleneck and
increasing speedup toward 2.0$\times$.

\subsection{Comparison with Production Systems}

Our single-node architecture differs fundamentally from production L2
systems (Table~\ref{tab:architecture_comparison}) in scale target:
Polygon, zkSync, and Scroll serve millions of users with horizontally
scaled prover pools, while our node serves a single enterprise. This
simplification eliminates the need for prover coordination, task
scheduling, and consensus mechanisms, but preserves the core
orchestration challenge: managing the asynchronous proof lifecycle while
maintaining continuous availability.

The push0 comparison is most informative. Our batch formation latency
(11.4~ms) is 2.2$\times$ push0's P50 (5.2~ms), attributable to our
node performing actual state tree operations (8 hash computations) versus
push0's pure message routing. When comparing equivalent operations
(message routing vs. batch formation without SMT), the overhead is
comparable.

\subsection{Limitations}

\begin{enumerate}
\item \textbf{Mock components}: Stage~1 uses calibrated mocks, not real
  components. Real SMT, snarkjs, and ethers.js may introduce variability
  not captured by deterministic mocks. Stage~2 will replace all mocks.

\item \textbf{Single enterprise}: All benchmarks use a single enterprise
  identifier. Concurrent enterprises sharing a node may introduce
  contention on the SMT and queue.

\item \textbf{Crash recovery not tested}: The WAL + checkpoint design
  is described but not empirically validated. Stage~3 will inject crashes
  at each pipeline phase.

\item \textbf{No network latency}: L1 submission uses a 2~s mock.
  Actual Avalanche Fuji latency may vary with network conditions.

\item \textbf{Pipeline simulation}: The pipeline comparison uses
  analytical simulation, not concurrent thread execution. Real pipelining
  may encounter event loop contention in Node.js.
\end{enumerate}
